% Generated draft: requires mathematical and source review.
\documentclass[11pt]{article}
\usepackage[margin=25mm]{geometry}
\usepackage{fontspec}
\setmainfont{Latin Modern Roman}
\usepackage{amsmath,amssymb,mathtools,unicode-math}
\setmathfont{Latin Modern Math}
\usepackage{enumitem,needspace,xurl,hyperref}
\setlength{\parindent}{0pt}
\setlength{\parskip}{5pt}
\setlength{\emergencystretch}{3em}
\newcommand{\N}{\mathbb{N}}
\newcommand{\Z}{\mathbb{Z}}
\newcommand{\Q}{\mathbb{Q}}
\newcommand{\R}{\mathbb{R}}
\newcommand{\C}{\mathbb{C}}
\newcommand{\F}{\mathbb{F}}
\newcommand{\A}{\mathbb{A}}
\newcommand{\PP}{\mathbb{P}}
\newcommand{\E}{\mathbb{E}}
\newcommand{\eps}{\varepsilon}
\newcommand{\dd}{\,\mathrm{d}}
\begin{document}
\section*{UnsolvedMath ID: 800330}
\textbf{Source problem code:} catalogue entry
\subsection{Definitions and mathematical statement}
For each integer \(n\geq 1\), let \(H_n\) be the hexagonal region in the unit triangular lattice whose interior angles are all \(120^\circ\) and whose consecutive side lengths are
\[
2n-1,\quad 2n,\quad 2n-1,\quad 2n-1,\quad 2n,\quad 2n-1.
\]
This uses the conventional interpretation of the source's ``\(2n-1,2n,2n-1\) semiregular hexagon.'' The source does not explicitly specify the domain of \(n\); positive integers are used here so that all side lengths are positive lattice lengths.

A lozenge is the union of two unit equilateral triangles sharing an edge. A lozenge tiling of \(H_n\) covers the region without gaps or overlapping interiors. Let \(\mathcal T_n\) be the finite set of its lozenge tilings. Let \(c_n\) be the triangular-lattice edge whose midpoint is the center of symmetry of \(H_n\), and let \(L_n\) be the lozenge formed by the two unit triangles incident to \(c_n\). These are the source's two innermost triangles.

The requested statement is that, for every integer \(n\geq 1\),
\[
3\bigl|\{T\in\mathcal T_n:L_n\in T\}\bigr|=|\mathcal T_n|.
\]
Equivalently, define \(G_n\) to have one vertex for each unit triangle of \(H_n\), with two vertices adjacent precisely when their triangles share an edge. Let \(e_n\) be the edge of \(G_n\) joining the vertices corresponding to the two triangles of \(L_n\). A perfect matching is a set of edges incident to every vertex exactly once. If \(M\) is chosen uniformly from the perfect matchings of \(G_n\), then the equivalent requested statement is
\[
\Pr(e_n\in M)=\frac13.
\]
\subsection{Short English statement}
For every positive integer \(n\), exactly one third of the lozenge tilings of the specified hexagon contain the lozenge at its center. Equivalently, a uniformly chosen perfect matching of the dual graph contains the central edge with probability \(1/3\). The positive-integer domain and the conventional hexagon geometry make explicit assumptions implicit in the supplied statement.
\subsection{Sources}
\begin{itemize}
\item Local UnsolvedMath registry: permanent ID \(800330\), exact source problem code \texttt{catalogue entry}, title ``Enumeration of Matchings — Problem 1,'' in the AMR Open Problem Lists collection. The supplied record attributes the problem to Propp, with proposed year 1998.
\item Propp, \emph{Twenty Open Problems in Enumeration of Matchings} (1998), Problem 1 (\#1). Supplied source URL: \url{https://arxiv.org/abs/math/9801060}.
\item The supplied provenance records access on 2026-10-16 and extraction by \texttt{arxiv\_latex}. These are supplied dated provenance claims, not independently verified access or source verification in this draft. The registry record was updated on 2026-10-16 and labels the problem ``open''; this is supplied curation, not an independently verified assessment of its mathematical status.
\end{itemize}
\end{document}
